\section{Further research questions and a concrete program}
\label{sec:research}

The proved reductions substantially narrow the relevant unknowns. The
next work should begin with the complementary parity components and
the next series-depth layer, while making every computational claim
replayable. The questions below distinguish mathematical obstructions,
finite relation computations, and implementation tasks. No proposed
numerical independence statement is needed to make progress on them.

\subsection{Complementary parity at levels three and four}

For the Gaussian single-endpoint doubles define
\[
 h_{a,b}=
 \begin{cases}
   \Re D_{a,b}(i),&a+b\text{ even},\\
   \Im D_{a,b}(i),&a+b\text{ odd}.
 \end{cases}
\]
The inversion formula used in this article does not determine these
components. Which of the entries through weights five to eight reduce
to the manuscript's single and half-point values after \emph{all
specified} regularized shuffle, stuffle, distribution, conjugation,
and parity relations are imposed? Which require additional transformed
arguments in an explicit functional reduction?

A concrete first deliverable is a sequence of exact rational matrices
with declared columns and rows. Include the product coordinates that
occur in the reductions, or quotient by a precisely defined span of
such products. If $V$ is the formal space and $R$ its proved relation
subspace, evaluation induces a map from $V/R$ onto the numerical span
under consideration. Therefore $\dim(V/R)$ provides an upper bound;
it does not automatically equal the numerical dimension. Each
eliminated coordinate should come with a row combination expressing
its reduction. This task can settle additional source identities even
when numerical independence remains out of reach.

\subsection{Relations beyond a selected double-shuffle presentation}

The mixed-root corrections show that a residual of one convergent
relation system can disappear under an already known functional
equation. The next comparison should measure this effect explicitly.
At each chosen weight and depth, start with a declared standard
relation module and adjoin the functional equations induced by
symmetries of the punctured projective line, together with their
endpoint corrections. Determine the exact increase in rank and
exhibit representatives of the added row-space directions.

At level four the relevant marked points are
$\{0,\infty,1,-1,i,-i\}$. One can calculate the fractional linear
transformations preserving this set and transport the iterated
integrals, keeping the basepoint and branch constants explicit. The
analogous level-three calculation should use its own marked points
and distribution relations, rather than assuming that the two
levels have identical relation counts. A comparison with the
constructive parity formulas of \cite{Panzer2017,Umezawa2025} should
be part of the relation specification. The objective is a verified
comparison of finite presentations; completeness for all numerical
relations is a separate question.

\subsection{The next depth layer: one index three or two indices two}

Theorem~\ref{one2:thm:closed} handles precisely one zero letter in the
positive-form binary word. At weight $N$ and depth $N-2$, the positive
indices necessarily consist of either one $3$ and otherwise $1$'s,
or two $2$'s and otherwise $1$'s. The corresponding words have two
zero letters. This is the natural next family, with a sharply
defined combinatorial scope.

Seek an explicit formula indexed by the numbers of $1$'s before,
between, and after the two zeros. A useful target is a finite
combination of classical and double polylogarithms at transformed
arguments, with factorial or binomial coefficients and all boundary
constants fixed. The one-zero proof suggests using a two-variable
logarithmic moment and then integrating one variable at a time.
At low weights, direct symbolic differentiation can verify candidate
functional formulas. At higher weights, an obstruction at the
function or coaction level would be more informative than a failed
fit to a chosen classical-value basket. Such an obstruction must
still be distinguished from numerical nonreducibility at a
particular algebraic argument.

\subsection{What survives modulo products at the half-Gaussian point?}

Corollary~\ref{one2:cor:primitive} proves that every imaginary one-$2$
value of weight $N$ is a signed binomial multiple of $\lambda_N$
modulo the explicitly stated lower-weight products. It is natural to
ask whether this quotient class survives, how it compares with the
other argument families, and whether the same binomial structure
has a useful interpretation in an algebra of motivic iterated
integrals.

The immediate task is to specify a motivic algebra and lift the
functional identity, rather than assigning a motivic meaning to a
formal numerical symbol. One can then calculate the coaction of
the remaining half-point class, retaining primitive period terms
that an ordinary symbol would discard. Proving that the image in
a particular motivic quotient is one-dimensional would strengthen
the current upper bound in that setting. Inferring numerical
independence would require an additional statement about the
period map and must be labeled accordingly. Even without such
an inference, the calculation could identify the correct candidate
generators for the next exact relation computation.

\subsection{Optimal coefficients and integer relation lattices}

The all-weight Bernoulli formula gives rational coefficient vectors
without numerical discovery. What are their sharp denominator and
divisibility patterns? Which primes occur in the least common
denominator at a given weight, and which apparent Bernoulli
denominators cancel because of the binomial factors? The
$\beta(w)$ coefficient and the absence of that term in $g_{1,w-1}$
are simple examples of structural information available before
any evaluation.

These questions can be pursued by exact symbolic arithmetic and
number-theoretic analysis. Publish the primitive integer rows,
their heights, and the integer lattice they generate. Hermite or
Smith normal forms can distinguish an integer lattice from its
saturation over $\mathbb Q$. A useful outcome would be a proved
height bound for the parity reductions and a principled choice of
coordinates for discovery searches. In particular, the low heights
of the two mixed real reductions in this article give concrete
tests that any replacement search implementation should recover.

\subsection{Exact certificates beyond the two-disk alphabet}

Corollary~\ref{cor:quadratic-certificate} already extends the proof
to a fixed imaginary quadratic field. Implementing it over
$\mathbb Q(\sqrt{-3})$ is a direct engineering task and would give
rational-radius certificates for the level-three and level-six
tables. This should precede a more ambitious general cyclotomic
implementation.

At larger conductors, some roots of unity are too close to $1$ for
the present reflected unit-disk majorant. For example, the distance
of $e^{2\pi i/q}$ from $1$ is $2\sin(\pi/q)$, which falls below $1$
when $q>6$. Develop rational splitting, path subdivision, or
transformations with a proved error bound that controls this
geometry. The desired theorem should display its dependence on
weight, precision, conductor, and coefficient-field degree. Shared
subwords should be reused so that subdividing paths does not
silently create an exponential expansion. A certificate should
remain replayable through an independent finite formula.

\subsection{Certified outcomes for bounded relation searches}

A search can be made mathematically informative without claiming
unbounded independence. Fix a finite constant basket and a height
bound $H$. For rational centers $r_j$ and certified errors
$\varepsilon_j$, every integer vector $c$ satisfying
$\max_j|c_j|\le H$ obeys
\[
 \left|\sum_jc_j\alpha_j-\sum_jc_jr_j\right|
 \le\sum_j|c_j|\varepsilon_j.
\]
Thus a finite exact lattice or enumeration certificate establishing
$|\sum_jc_jr_j|>\sum_j|c_j|\varepsilon_j$ for every permitted
nonzero $c$ excludes all relations within that declared range.
Already proved dependencies must first be handled with the
height constraints translated correctly to the remaining
coordinates.

A robust program should have three outcomes: a candidate upgraded
to a proved identity by an exact functional-equation or relation-row
certificate; certified exclusion within the finite bound; or an
unresolved outcome. An enclosure containing zero belongs to the
last category until an identity proof is supplied. Such a program
would make the source's failed searches reproducible and clearly
state what their failure actually establishes.

\subsection{Recommended sequence of deliverables}

The most economical next stage is to integrate the eight proof
upgrades and the four mixed corrections, then publish the exact
relation matrices through weight six with their elimination
certificates. Extend those matrices to weights seven and eight
only after the earlier rows are reproducible. In parallel, derive
the two-zero moment formula and implement the quadratic-field
certificate specialization. The deeper questions about motivic
classes and numerical independence can then be asked of a
well-defined collection of remaining values, with the known
parity reductions already removed.

The proved results here supply an infinite family of exact
identities, rather than a finite list of guesses: the even-weight
Gaussian double rule, every position of the one-$2$ family, the
inverse-color parity specialization, and the uniform certificate
theorem all admit systematic extension. Their explicit hypotheses
and executable coefficient formulas provide concrete starting
points for that next research stage.
